% =============================================================================
\section{Modélisation et approche statique}
% =============================================================================

% Ensembles, paramètres et variables.
\begin{frame}{Ce qu'il faut décider : variables et paramètres}
  \small
  \begin{columns}[T,onlytextwidth]
    \begin{column}{0.52\textwidth}
      \begin{block}{Ensembles et décision}
        \begin{tabular}{@{}ll@{}}
          $P = \{1,\dots,n\}$ & patients à planifier \\
          $S = \{1,\dots,m\}$ & salles opératoires \\
          $T = \{1,\dots,H\}$ & jours de l'horizon \\
          $x_{i,s,t} \in \{0,1\}$ & patient $i$ en salle $s$ le jour $t$ \\
        \end{tabular}
        \vspace{4pt}
        \[
          \sum_{s \in S}\sum_{t \in T} x_{i,s,t} = 1
          \quad \forall i \in P
        \]
      \end{block}
    \end{column}
    \begin{column}{0.46\textwidth}
      \begin{block}{Paramètres issus des données}
        \begin{tabular}{@{}ll@{}}
          $d_i$ & durée opératoire du patient $i$ \\
          $S_i$ & durée de séjour (jours) \\
          $a_i$ & ambulatoire ($1$) ou non ($0$) \\
          $C_{\text{vac}}$ & capacité d'une vacation \\
          $L_{\max}$ & capacité en lits \\
          $A_{\max}$ & capacité en places ambu. \\
          $\tau_{\text{Urg}}(t)$ & marge urgences du jour $t$ \\
        \end{tabular}
      \end{block}
    \end{column}
  \end{columns}
  \source{Durées $d_i$ et $S_i$ estimées à partir de l'historique ; capacité de
    vacation typique : 4\,h (240 min).}
\end{frame}

% Fonction objectif : trois tensions.
\begin{frame}{Fonction objectif : trois tensions à équilibrer}
  \footnotesize
  \[
    \begin{aligned}
      \min\; & \alpha \sum_{J \in T}\bigl(L(J)-L_{\text{cible}}(J)\bigr)^2
             + \beta \sum_{J \in T}\bigl(A(J)-A_{\text{cible}}(J)\bigr)^2 \\
             & + \gamma \sum_{J \in T}\Delta_{\text{bloc}}(J)^2
    \end{aligned}
  \]
  \begin{columns}[T,onlytextwidth]
    \begin{column}{0.52\textwidth}
      \begin{block}{Lecture}
        \begin{itemize}
          \item terme 1 : pénalise les \textbf{pics et creux} de lits ;
          \item terme 2 : lisse la charge des \textbf{places ambulatoires} ;
          \item terme 3 : pénalise l'écart à la capacité du bloc, urgences
            comprises ;
          \item $\alpha,\beta,\gamma>0$ : poids à régler.
        \end{itemize}
      \end{block}
    \end{column}
    \begin{column}{0.46\textwidth}
      \begin{alertblock}{Coût réellement implémenté}
        \footnotesize
        \[
          \begin{aligned}
            \text{coût} &= 5 \cdot \text{dépassement}_{\text{vac}}
                          + 3 \cdot \text{dépassement}_{\text{lits}} \\
                         &\quad + 0{,}05 \cdot \text{écart-type}_{\text{vac}}
          \end{aligned}
        \]
        La \textbf{fitness} vaut $-\text{coût}$ : on maximise, $0$ = planning
        parfaitement faisable et équilibré.
      \end{alertblock}
    \end{column}
  \end{columns}
  \source{Le terme bloc et le terme lits sont des contraintes dures relâchées ;
    l'ambulatoire intervient dans la formulation cible et via le type de prise
    en charge prédit.}
\end{frame}

% Les trois métaheuristiques d'origine et leur rôle.
\begin{frame}{Trois métaheuristiques complémentaires}
  \footnotesize
  Le problème étant NP-difficile, l'optimum exact n'est pas atteignable à
  l'échelle d'un hôpital : on cherche une bonne solution en un temps borné.

  \vspace{0.15cm}
  \renewcommand{\arraystretch}{1.12}
  \begin{tabular}{p{3.2cm}p{4.3cm}p{4.6cm}}
    \toprule
    \textbf{Méthode} & \textbf{Rôle dans le projet} & \textbf{Force} \\
    \midrule
    Recuit simulé & Construction initiale : accepte parfois de « redescendre »
      & Large exploration, évite les optima locaux \\
    Recherche Tabou & Raffinement local : meilleur voisin non tabou
      & Mémorise les mouvements, évite le cyclage \\
    Algorithme génétique & Recherche de compromis global : population
      & Diversité, combine de bonnes briques \\
    \bottomrule
  \end{tabular}

  \vspace{0.2cm}
  \begin{alertblock}{À retenir}
    Trois logiques \textbf{complémentaires --- pas rivales} : explorer,
    intensifier, diversifier.
  \end{alertblock}
  \source{Chaque méthode partage la même interface
    \texttt{solve(instance) -> planning} pour une comparaison à conditions
    égales.}
\end{frame}

% Validation exacte puis limite du statique.
\begin{frame}{Validation par force brute et limite du statique}
  \small
  \begin{columns}[T,onlytextwidth]
    \begin{column}{0.5\textwidth}
      \begin{exampleblock}{Preuve de correction}
        \footnotesize Sur l'instance jouet (7 patients, 4 vacations,
        $4^7 = 16\,384$ combinaisons), l'optimum est calculé par énumération
        complète : \textbf{les 10 méthodes l'atteignent sur les 20 graines}.
      \end{exampleblock}
      \centering
      \includegraphics[width=0.92\linewidth,height=0.26\textheight,keepaspectratio]{petite_taux_optimum.png}
    \end{column}
    \begin{column}{0.48\textwidth}
      \begin{alertblock}{La limite}
        \footnotesize Le planning obtenu est optimal \textbf{pour les durées
        prévues}. Dès qu'une urgence arrive, qu'un patient annule, qu'un lit
        ferme ou qu'un bloc prend du retard, l'optimum statique n'est plus le
        bon planning.
      \end{alertblock}
      \begin{block}{Conséquence}
        \footnotesize Il faut passer d'un planning \textbf{optimal} à un
        planning \textbf{robuste}, capable d'absorber les aléas.
      \end{block}
    \end{column}
  \end{columns}
  \source{Instance de validation à 7 patients ; l'optimum exact sert de test de
    correction, pas de départage entre méthodes.}
\end{frame}

